\documentclass[11pt]{article}
\usepackage[utf8]{inputenc}
\usepackage[T1]{fontenc}
\usepackage{amsmath,amssymb,amsfonts,graphicx,color,xcolor,physics,bm,geometry,siunitx}
\usepackage[numbers,super,sort&compress]{natbib}
\usepackage{xr-hyper}
\usepackage{hyperref}
\usepackage{cleveref}
\externaldocument{Manuscript_JPCL_26-06-13}
\externaldocument[SI-]{SI_JPCL_26-06-13}
\geometry{margin=2.cm}

\hypersetup{
    colorlinks=true,
    linkcolor=blue,
    citecolor=blue,
    urlcolor=blue
}

\begin{document}

\begin{flushright}
    Douala, June 17, 2026
\end{flushright}

\noindent \textbf{Prof. Gregory Scholes}\\
Editor-in-Chief\\
\textit{The Journal of Physical Chemistry Letters}\\
\textit{ACS Publications}

\vspace{1em}

\noindent \textbf{Re: Manuscript ID jz-2026-00994t.R1 (Major Revision)}\\
\textbf{Title:} Selective Vibronic Excitation for Coherent Energy Transport in Photosynthetic and Agrivoltaic Systems

\vspace{2em}

\noindent We thank the reviewers for their continued insightful and constructive evaluation of our revised manuscript. We have carefully addressed the remaining concerns raised by Reviewer 2 and Reviewer 3. In particular, we have rigorously purged the manuscript of all software engineering jargon, clarified the theoretical definition of the transfer yield, and explicitly connected our discussion to established vibronic basis formalisms.

Below, we provide a point-by-point response to the new comments. The corresponding changes in the revised manuscript and Supporting Information (SI) are highlighted.

\vspace{2em}

\section*{Response to Reviewers}

\subsection*{Reviewer 2}

\textbf{(i) Connection to early work on vibronic and fractional resonance}
\textit{Comment: As stated in the previous review, vibronic resonance and fractional resonance have been analyzed extensively [e.g., JPC. Lett. 6(4), 627 (2015); 13, 6831 (2022)], and the authors should connect to the early work.}

\textbf{Response:} These references are now integrated into the Introduction and Discussion. We explicitly link our selective excitation scheme to the resonance mechanisms described in \textit{J. Phys. Chem. Lett.} 6, 627 (2015) and \textit{J. Phys. Chem. Lett.} 13, 6831 (2022), framing our active spectral control within the established consensus on resonance-enhanced transport.

\subsection*{Reviewer 3}

\textbf{1. Definition of the Forward Transfer Yield}
\textit{Comment: The definition for the forward transfer yield is basically 1 minus the survival probability... transfer yield should be defined as the long-time (equilibrium) population of that state...}

\textbf{Response:} We agree that 1 minus the survival probability of the initial state is an inadequate metric for transfer efficiency. We have redefined the forward transfer yield ($\Phi_{\mathrm{FT}}$) as the long-time equilibrium population of the target state (BChl 3). The corrected definition is now provided in Eq. 7. All associated data and Table 1 have been updated accordingly.

\textbf{2. Spectral Density and "189 Modes"}
\textit{Comment: I still don’t see the spectral density in the paper... I don’t understand how that number [189 modes] was obtained.}

\textbf{Response:} For the 7-site FMO model, coupling each site to 12 localized vibronic modes yields 84 total modes. The mention of "189 modes" was an error from an exploratory model using inter-site correlated baths, which was not used in the production runs. This is corrected. The total reorganization energy ($\lambda_{\mathrm{total}} \approx \SI{250}{\per\centi\meter}$) is now explicitly broken down in SI Table S1.

\textbf{3. Raw Simulation Results and Decoherence}
\textit{Comment: The raw simulation results are puzzling... Fig. 3b does not show a decay.}

\textbf{Response:} The high-frequency oscillations previously visualized masked the long-time limit. We have updated our ensemble simulation to extend the propagation time. Revised figures now display the long-time envelope decay of the $l_1$-norm coherence and site population equilibration, confirming environmentally-induced decoherence.

\textbf{4. PT-HOPS Simulation Method and Scaling}
\textit{Comment: Does this imply that the cost of a 2-state calculation is the same as the cost of a 100-state calculation? ... How can just 2 Matsubara terms lead to converged results?}

\textbf{Response:} We have revised the manuscript to remove overstatements regarding PT-HOPS scaling. We clarify that while SBD effectively mitigates the exponential scaling with respect to bath modes, the computational cost remains polynomial with system size. Regarding Matsubara convergence, since the Drude cutoff $\gamma_D = \SI{50}{\per\centi\meter}$ is below the first Matsubara frequency $\nu_1 \approx \SI{205}{\per\centi\meter}$ at room temperature, the expansion converges rapidly at $K=2$. This is mathematically justified in SI Section S2.

\vspace{3em}

\noindent We have also explicitly verified the convergence of the hierarchy depth (L) at L=8, demonstrating an MAE of 3e-5 relative to L=7, ensuring the robustness of our results.

We believe that these extensive revisions have resolved the critical issues identified, ensuring the scientific rigor and presentation quality required by \textit{The Journal of Physical Chemistry Letters}.

\vspace{1.5em}

\noindent Sincerely,

\vspace{2.5em}

\noindent \textbf{Steve Cabrel Teguia Kouam}\\
Corresponding Author\\
University of Douala, Cameroon

\end{document}
